\documentclass[a4paper,11pt]{article}
\usepackage[T2A]{fontenc}
\usepackage[utf8]{inputenc}
\usepackage[russian]{babel}
\usepackage{amsmath,amssymb,mathtools}
\usepackage{helvet}
\renewcommand{\familydefault}{\sfdefault}
\usepackage{geometry}
\geometry{left=22mm,right=20mm,top=22mm,bottom=23mm}
\usepackage{microtype}
\usepackage{icomma}
\usepackage{booktabs,tabularx,array}
\usepackage{enumitem}
\usepackage{xcolor}
\usepackage{tikz}
\usepackage{hyperref}
\usepackage{fancyhdr}
\usepackage{setspace}
\setstretch{1.14}
\definecolor{accent}{HTML}{006F75}
\definecolor{darkaccent}{HTML}{00464D}
\definecolor{ink}{HTML}{25343A}
\definecolor{err}{HTML}{AD2C32}
\hypersetup{colorlinks=true,linkcolor=darkaccent,pdfauthor={Лёвин Артём Александрович},pdftitle={Ревью домашней работы — Ярослав Гаврилов — 16.11.2026}}
\setlength{\parindent}{0pt}
\setlength{\parskip}{3pt}
\setlist[itemize]{nosep,topsep=3pt,leftmargin=1.5em}
\pagestyle{fancy}
\fancyhf{}
\fancyhead[L]{\small\color{darkaccent}Ревью домашней работы}
\fancyhead[R]{\small\color{ink}16.11.2026}
\fancyfoot[C]{\small\color{ink}\thepage}
\renewcommand{\headrulewidth}{0.25pt}
\setlength{\headheight}{15pt}
\newcommand{\rlabel}[1]{\par\noindent\textbf{\color{darkaccent}#1}\quad}
\newcommand{\errorlabel}{\textcolor{err}{\textbf{Ошибка / недочёт:}}}
\newcommand{\keylabel}{\textcolor{blue!60!black}{\textbf{Ключевая идея:}}}
\newcommand{\finallabel}{\textcolor{teal!60!black}{\textbf{Итог:}}}
\newcommand{\tasktitle}[2]{\clearpage\phantomsection\label{task:#1}\section*{Задание №\,#1. #2}\vspace{-3mm}\textcolor{accent}{\rule{\linewidth}{0.45pt}}}
\begin{document}
\thispagestyle{empty}
\begin{center}
\vspace*{30mm}
{\Large\bfseries\color{darkaccent}МАТЕМАТИКА}\par
\vspace{24mm}
{\huge\bfseries Ревью домашней работы}\par
\vspace{6mm}
{\Large Ярослав Гаврилов}\par
\vspace{10mm}
{\large Задачи на движение по прямой}\par
\vspace{12mm}
{\large 16 ноября 2026 года}\par
\vfill
{\large Лёвин Артём Александрович}\par
{\small эксклюзивно для Ярослава Гаврилова}\par
\vspace{12mm}
\begin{tikzpicture}
\draw[accent!65,line width=0.9pt] (-5,0) .. controls (-2.6,0.7) and (-1.1,-0.55) .. (0,0) .. controls (1.6,0.7) and (3.0,-0.5) .. (5,0);
\end{tikzpicture}
\end{center}
\clearpage
\section*{Сводная таблица}
В таблице находятся только задания, которые нужно исправить или завершить. Нажмите на номер для перехода к разбору.
\vspace{4mm}\par
\noindent\begingroup\small
\setlength{\tabcolsep}{1.8pt}
\renewcommand{\arraystretch}{1.35}
\begin{tabularx}{\linewidth}{@{}>{\centering\arraybackslash}p{9mm} >{\raggedright\arraybackslash}p{30mm} X >{\raggedright\arraybackslash}p{28mm} >{\raggedright\arraybackslash}p{35mm}@{}}
\toprule
\textbf{№}&\textbf{Тема}&\textbf{Суть ошибки / недочёта}&\textbf{Ваш ответ}&\textbf{Верный ответ}\\
\midrule
\hyperref[task:7]{7}&Движение с поздним стартом&При раскрытии знаменателя пропущен множитель $x$; решение не завершено.&Не указан&Грузовик --- $60$, автомобиль --- $80$ км/ч.\\
\hyperref[task:8]{8}&Догонка поездов&После номера задания решение отсутствует.&Не указан&$90$ км от станции.\\
\hyperref[task:9]{9}&Движение с разворотом&Не представлены две таблицы и система уравнений.&Не представлен&Система (см. разбор).\\
\hyperref[task:10]{10}&Догонка и общий финиш&Не представлены таблицы и два требуемых уравнения.&Не представлен&Система (см. разбор).\\
\bottomrule
\end{tabularx}
\endgroup
\clearpage
\thispagestyle{empty}
\vspace*{0.37\textheight}
\begin{center}
{\Huge\bfseries\color{darkaccent}Разбор ошибок}
\end{center}

\tasktitle{7}{Одинаковое расстояние, разные моменты старта}
\rlabel{Текст задания.} Грузовик и автомобиль проезжают путь $60$ км каждый. Грузовик выехал первым; автомобиль отправился через $15$ минут, движется на $20$ км/ч быстрее и прибывает одновременно с грузовиком. Найдите скорость грузовика.

\rlabel{Ваше решение.} Вы обозначили скорость грузовика через $x$ км/ч, скорость автомобиля --- через $x+20$ км/ч, а время движения --- соответственно $60/x$ и $60/(x+20)$ часов. Затем получили:
\[
\frac{60}{x}-\frac{60}{x+20}=\frac{15}{60}=\frac14,
\]
\[
\frac{60(x+20)-60x}{x(x+20)}=\frac14.
\]
Далее в записи знаменатель превращается в $x^2+20$, после чего преобразование остаётся незавершённым.

\rlabel{Ваш ответ.} Не указан.

\par\noindent\errorlabel\ Пропущен множитель $x$ при раскрытии скобок в знаменателе. Кроме того, исходная дробь далее преобразуется без корректного умножения обеих частей уравнения на общий знаменатель.

\rlabel{Где именно возникает ошибка.} Последний полностью правильный шаг:
\[
\frac{60(x+20)-60x}{x(x+20)}=\frac14.
\]
Следующий шаг должен быть
\[
 x(x+20)=x^2+20x,
\]
а не $x^2+20$. В записи $x^2+20$ потерян множитель $x$ при произведении $x\cdot20$.

\rlabel{Почему это неверно.} Распределительный закон требует умножить $x$ на \emph{каждое} слагаемое в скобках:
\[
 x\cdot(x+20)=x\cdot x+x\cdot20=x^2+20x.
\]
Простая проверка: при $x=10$ исходное произведение равно $10\cdot30=300$, тогда как ошибочное выражение $10^2+20=120$. Они не равны. Поэтому дальнейшие действия с ошибочным знаменателем уже не сохраняют исходное уравнение.

\par\noindent\keylabel\ Разность времён для одинакового расстояния равна задержке старта. После приведения дробей к общему знаменателю удобнее выполнить перекрёстное умножение и решить квадратное уравнение, обязательно проверив положительность скорости.

\rlabel{Исправление от последнего правильного шага.} В числителе раскрываем скобки, а знаменатель оставляем произведением:
\[
\frac{60x+1200-60x}{x(x+20)}=\frac14,
\qquad \frac{1200}{x(x+20)}=\frac14.
\]
Поскольку $x>0$, знаменатель не равен нулю. Умножаем обе части на $4x(x+20)$:
\[
 4800=x(x+20)=x^2+20x.
\]

\rlabel{Полное правильное решение.} Пусть $x$ км/ч --- скорость грузовика. Тогда автомобиль едет со скоростью $x+20$ км/ч. Обе скорости положительны, то есть $x>0$. Формула времени:
\[
 t=\frac{s}{v}.
\]
Получаем таблицу движения:
\[
\begin{array}{lccc}\toprule
 & v\ (\text{км/ч}) & t\ (\text{ч}) & s\ (\text{км})\\\midrule
 \text{Грузовик} & x & \frac{60}{x} & 60\\
 \text{Автомобиль} & x+20 & \frac{60}{x+20} & 60\\\bottomrule
\end{array}
\]
Автомобиль стартовал на четверть часа позже и прибыл одновременно, следовательно, его время движения на четверть часа меньше:
\[
\begin{aligned}
\frac{60}{x}-\frac{60}{x+20}&=\frac14,\\
\frac{60(x+20)-60x}{x(x+20)}&=\frac14,\\
\frac{1200}{x(x+20)}&=\frac14,\\
x^2+20x-4800&=0.
\end{aligned}
\]
Разложим левую часть на множители:
\[
 x^2+20x-4800=(x-60)(x+80).
\]
Поэтому $x=60$ или $x=-80$. Отрицательная скорость не подходит по смыслу задачи. Значит, скорость грузовика равна $60$ км/ч, а автомобиля --- $80$ км/ч.

\rlabel{Проверка.} Грузовик проезжает $60$ км за $60/60=1$ час. Автомобиль --- за $60/80=3/4$ часа, то есть за $45$ минут. Разность времён равна $15$ минут. Условие выполнено.

\par\noindent\finallabel\ Грузовик движется со скоростью $60$ км/ч; автомобиль --- $80$ км/ч.

\rlabel{Задача на закрепление.} Грузовик проехал $72$ км. Легковой автомобиль отправился по тому же маршруту на $30$ минут позже, двигался на $12$ км/ч быстрее и прибыл одновременно с грузовиком. Найдите скорости обоих автомобилей.

\tasktitle{8}{Поезда: задержка отправления и догонка}
\rlabel{Текст задания.} Грузовой поезд выехал со станции со скоростью $54$ км/ч. Через $40$ минут в том же направлении отправился скорый поезд со скоростью $90$ км/ч. Найдите расстояние от станции до места, где скорый поезд догонит грузовой. Требуется составить уравнение разности времён.

\rlabel{Ваше решение.} На фотографии присутствует обозначение задания №\,8, но расчёты и ответ отсутствуют.

\rlabel{Ваш ответ.} Не указан.

\par\noindent\errorlabel\ Задание не завершено. В частности, не составлено требуемое уравнение разности времён.

\rlabel{Где именно работа остановилась.} Последнего записанного математического шага для этого задания нет. Первый необходимый шаг --- перевести $40$ минут в часы:
\[
40\ \text{мин}=\frac{40}{60}\ \text{ч}=\frac23\ \text{ч}.
\]

\rlabel{Почему это важно.} Расстояние одно и то же для обоих поездов: от станции до точки встречи. Но грузовой поезд движется дольше, поскольку выехал раньше. Уравнение должно сравнивать \emph{времена}, а не складывать скорости или пройденные пути.

\par\noindent\keylabel\ Если оба участника проходят до встречи одинаковое расстояние $x$, то их времена равны $x/v_1$ и $x/v_2$. Разность времен равна задержке более позднего старта.

\rlabel{Исправление: как начать решение.} Обозначим через $x$ км расстояние до места встречи. Времена движения составляют $x/54$ и $x/90$ часа. Грузовой поезд ехал на $2/3$ часа дольше:
\[
\frac{x}{54}-\frac{x}{90}=\frac23.
\]

\rlabel{Полное правильное решение.} Заполним таблицу:
\[
\begin{array}{lccc}\toprule
 & v\ (\text{км/ч}) & t\ (\text{ч}) & s\ (\text{км})\\\midrule
\text{Грузовой} &54&\frac{x}{54}&x\\
\text{Скорый} &90&\frac{x}{90}&x\\\bottomrule
\end{array}
\]
Отсюда:
\[
\begin{aligned}
\frac{x}{54}-\frac{x}{90}&=\frac23,\\
\frac{5x-3x}{270}&=\frac23,\\
\frac{2x}{270}&=\frac23,\\
\frac{x}{135}&=\frac23,\\
x&=90.
\end{aligned}
\]
Вычисление можно проверить другим способом. За первые $40$ минут грузовой поезд создаёт отрыв $54\cdot\frac23=36$ км. Скорый поезд сокращает отрыв со скоростью $90-54=36$ км/ч, поэтому догоняет его через $36/36=1$ час после собственного отправления. За этот час скорый поезд проезжает $90$ км.

\rlabel{Проверка.} Грузовой поезд движется до встречи $90/54=5/3$ часа, или $1$ час $40$ минут. Скорый --- $90/90=1$ час. Разность ровно $40$ минут.

\par\noindent\finallabel\ Скорый поезд догонит грузовой на расстоянии $90$ км от станции.

\rlabel{Задача на закрепление.} Грузовой поезд отправился со скоростью $60$ км/ч. Через $30$ минут в том же направлении выехал скорый поезд со скоростью $100$ км/ч. Составьте уравнение разности времён и найдите расстояние от станции до места встречи.

\tasktitle{9}{Разворот мотоциклиста: две таблицы и система}
\rlabel{Текст задания.} Между городами $A$ и $B$ --- $120$ км. Автомобиль выехал из $A$ в $B$. Спустя $15$ минут оттуда же отправился мотоциклист со скоростью $80$ км/ч. Он догнал автомобиль в точке $C$ и сразу повернул обратно. Пока автомобиль доехал от $C$ до $B$, мотоциклист проехал назад треть расстояния $AC$. Обозначьте $AC=x$ км, скорость автомобиля --- $y$ км/ч. Требуется нарисовать схему, составить две таблицы и получить систему уравнений. Решать систему не нужно.

\rlabel{Ваше решение.} В представленных фотографиях решение задания №\,9 отсутствует.

\rlabel{Ваш ответ.} Система уравнений не представлена.

\par\noindent\errorlabel\ Отсутствуют требуемые этапы: разбиение движения на два промежутка, две таблицы и система.

\rlabel{Где именно работа остановилась.} Последнего записанного шага нет. Первое необходимое действие --- зафиксировать два \emph{разных} события: догонку в точке $C$ и последующий одновременный интервал движения автомобиля до $B$ и мотоциклиста обратно.

\rlabel{Почему важно разделить движение.} До встречи оба едут из $A$ в $C$ и проходят одинаковое расстояние $x$, но автомобиль начал на $15$ минут раньше. После встречи пути различны: автомобиль проходит $CB=120-x$, а мотоциклист --- $x/3$. Зато время движения на втором этапе одинаково. Если ошибочно приравнять расстояния второго этапа или сложить время первого этапа со вторым, получится система, не описывающая условие.

\par\noindent\keylabel\ Первый этап даёт \emph{разность} времён, второй --- их \emph{равенство}. Для каждого этапа используйте свою таблицу.

\rlabel{Исправление: схема движения.} Точка $C$ находится между $A$ и $B$:
\[
\underbrace{A\quad\longrightarrow\quad C}_{AC=x}
\quad\underbrace{\longrightarrow\quad B}_{CB=120-x}.
\]
До $C$ оба едут вправо; после $C$ автомобиль продолжает вправо к $B$, а мотоциклист едет влево, преодолевая $x/3$ км.

\rlabel{Полное правильное решение.} \textbf{Первая таблица --- до догонки.} Оба проезжают $AC=x$ км. Скорость автомобиля $y$ км/ч, скорость мотоциклиста $80$ км/ч:
\[
\begin{array}{lccc}\toprule
 & v\ (\text{км/ч}) & t\ (\text{ч}) & s\ (\text{км})\\\midrule
\text{Автомобиль} & y & \frac{x}{y} & x\\
\text{Мотоциклист} &80&\frac{x}{80}&x\\\bottomrule
\end{array}
\]
Автомобиль выехал раньше на $15/60=1/4$ часа, следовательно:
\[
\boxed{\frac{x}{y}-\frac{x}{80}=\frac14.}\tag{1}
\]
\textbf{Вторая таблица --- после догонки.} В момент встречи участники одновременно начинают второй этап:
\[
\begin{array}{lccc}\toprule
 & v\ (\text{км/ч}) & t\ (\text{ч}) & s\ (\text{км})\\\midrule
\text{Автомобиль, }C\to B & y & \frac{120-x}{y} & 120-x\\
\text{Мотоциклист, обратно} &80&\frac{x/3}{80}=\frac{x}{240}&\frac{x}{3}\\\bottomrule
\end{array}
\]
Указанные участки пройдены за одно и то же время. Поэтому:
\[
\boxed{\frac{120-x}{y}=\frac{x}{240}.}\tag{2}
\]
Объединяем полученные уравнения:
\[
\boxed{\begin{cases}
\dfrac{x}{y}-\dfrac{x}{80}=\dfrac14,\\[5pt]
\dfrac{120-x}{y}=\dfrac{x}{240}.
\end{cases}}
\]
По смыслу задачи $0<x<120$ и $0<y<80$: точка догонки внутренняя, а догоняющий должен быть быстрее автомобиля.

\rlabel{Проверка составления.} В первом уравнении обе дроби имеют размерность часов, и первая больше второй на четверть часа. Во втором уравнении слева и справа также записано время: расстояние до $B$, делённое на скорость автомобиля, равно расстоянию обратного хода, делённому на скорость мотоциклиста. Таким образом, обе части системы соответствуют двум разным этапам движения.

\par\noindent\finallabel\ Требуемая система получена; решать её по условию не нужно.

\rlabel{Задача на закрепление.} Между $A$ и $B$ --- $150$ км. Автомобиль выехал первым, а через $10$ минут вслед за ним отправился мотоциклист со скоростью $90$ км/ч. Догнав автомобиль в точке $C$, мотоциклист повернул обратно и за время движения автомобиля от $C$ до $B$ проехал назад четверть расстояния $AC$. Обозначьте $AC=x$ км, скорость автомобиля --- $y$ км/ч. Составьте две таблицы и систему; решать её не нужно.

\tasktitle{10}{Догонка и общий финиш: два этапа}
\rlabel{Текст задания.} Из $A$ в $B$ по дороге длиной $270$ км выехал грузовик. Через полчаса следом отправился автомобиль, скорость которого в три раза выше. Через некоторое время автомобиль догнал грузовик, и оба продолжили движение с прежними скоростями. Автомобиль прибыл в $B$ на $45$ минут раньше. Обозначьте скорости грузовика и автомобиля через $v$ и $w$. Требуется составить отдельные таблицы до догонки и после догонки и записать два уравнения, связывающих $v$ и $w$. Решать систему не требуется.

\rlabel{Ваше решение.} В представленных фотографиях решение задания №\,10 отсутствует.

\rlabel{Ваш ответ.} Не представлен.

\par\noindent\errorlabel\ Не составлены таблицы двух этапов движения и не записаны два требуемых соотношения.

\rlabel{Где именно работа остановилась.} Последнего записанного шага нет. Необходимо начать с точного определения скоростей:
\[
 v>0,\qquad w=3v.
\]
Затем следует учесть оба временных сдвига: отправление автомобиля на $1/2$ часа позже и его прибытие на $3/4$ часа раньше.

\rlabel{Почему здесь нельзя смешивать интервалы.} От старта грузовика до догонки он движется на полчаса дольше автомобиля; от догонки до города $B$ оба стартуют одновременно, но прибывают с разницей в $45$ минут. Если использовать всю длину $270$ км как путь после догонки, получится неверная таблица. Для второго этапа необходимо вычесть путь, пройденный до встречи.

\par\noindent\keylabel\ До догонки равны \emph{расстояния} от $A$ до точки встречи. После догонки различаются \emph{времена} прохождения общего оставшегося расстояния. Удобно сначала ввести вспомогательное расстояние $d$, а затем исключить его.

\rlabel{Исправление: модель первого этапа.} Пусть $d$ км --- расстояние от $A$ до места догонки. Тогда $0<d<270$. Грузовик шёл до встречи $d/v$ часа, а автомобиль --- $d/w$ часа. Значит,
\[
 \frac{d}{v}-\frac{d}{w}=\frac12.\tag{1}
\]

\rlabel{Полное правильное решение.} \textbf{Первая таблица --- до догонки.}
\[
\begin{array}{lccc}\toprule
 & v\ (\text{км/ч}) & t\ (\text{ч}) & s\ (\text{км})\\\midrule
\text{Грузовик} & v & \frac{d}{v}&d\\
\text{Автомобиль}&w&\frac{d}{w}&d\\\bottomrule
\end{array}
\]
Грузовик стартует на $1/2$ часа раньше. Получаем уравнение (1).

\textbf{Вторая таблица --- после догонки.} Оба находятся в одной точке и проезжают до $B$ расстояние $270-d$ км:
\[
\begin{array}{lccc}\toprule
 & v\ (\text{км/ч}) & t\ (\text{ч}) & s\ (\text{км})\\\midrule
\text{Грузовик} & v & \frac{270-d}{v}&270-d\\
\text{Автомобиль}&w&\frac{270-d}{w}&270-d\\\bottomrule
\end{array}
\]
После догонки автомобиль заканчивает путь на $45/60=3/4$ часа раньше, следовательно,
\[
\frac{270-d}{v}-\frac{270-d}{w}=\frac34.\tag{2}
\]
\textbf{Исключаем вспомогательное расстояние.} Раскрываем скобки во втором уравнении:
\[
270\left(\frac1v-\frac1w\right)-d\left(\frac1v-\frac1w\right)=\frac34.
\]
Из первого уравнения известно, что
\[
 d\left(\frac1v-\frac1w\right)=\frac12.
\]
Поэтому
\[
270\left(\frac1v-\frac1w\right)-\frac12=\frac34,
\]
\[
270\left(\frac1v-\frac1w\right)=\frac34+\frac12=\frac54.
\]
Поскольку автомобиль едет втрое быстрее грузовика, окончательно получаем два уравнения, связывающие только $v$ и $w$:
\[
\boxed{\begin{cases}
w=3v,\\[5pt]
\dfrac{270}{v}-\dfrac{270}{w}=\dfrac54.
\end{cases}}
\]
Решать систему по условию не требуется.

\rlabel{Проверка составления.} Разность полных \emph{времён движения} грузовика и автомобиля равна $1/2+3/4=5/4$ часа: грузовик начал раньше на полчаса и закончил позже на три четверти часа. Поэтому второе уравнение согласуется с хронологией задачи. Скорости положительны, $w=3v>v$; положение догонки должно удовлетворять $0<d<270$.

\par\noindent\finallabel\ Получены две таблицы по этапам и система из двух уравнений для скоростей $v$ и $w$.

\rlabel{Задача на закрепление.} Между пунктами $A$ и $B$ --- $360$ км. Грузовик отправился первым, а через $20$ минут за ним выехал автомобиль со скоростью вдвое выше. Автомобиль догнал грузовик и прибыл в $B$ на $50$ минут раньше него. Обозначьте скорости через $v$ и $w$. Составьте таблицы движения до и после догонки и запишите два уравнения для $v$ и $w$. Решать систему не нужно.

\end{document}